\section{Discussion}
\label{sec:discussion}

\paragraph{What the audit changes in the diagnosis.}
The earlier failure was attributed to a bound. The audit shows that, for a
verified projection, the bound is never reached. The actual obstruction lies
in the \emph{image} of the parameterisation: the FO-LQI synthesis produces
neither the low values of $k_p$ nor the values $b<1$ required by the feasible
points found by direct search. A bound is therefore meaningful only relative
to what the synthesis chain can produce; a saturation rate says nothing
without the distinction between structural and pathological saturation.

\paragraph{What the campaign establishes.}
Over 540 preregistered cells, no controller is robustly feasible. This
negative result is measured, not assumed: sterile seeds are reported, Wilson
intervals bound the success rate, and the sealed test remained untouched for
lack of a candidate to test. Freeing $b$ or $k_p$ makes \emph{nominal}
feasibility reachable, but not robustness: tuning freedom is not enough.

\paragraph{A physical obstruction.}
The voltage floor at minimum duty ratio is a property of the converter, not of
the controller. The threshold $V_\mathrm{min}\ge0.8$~V lies above this floor
in most uncertainty cases, so a control action saturated during the
$24\to1$~V descent cannot satisfy it. The specification thus implicitly
requires the controller to leave saturation in time, a requirement that
depends on anti-windup as much as on gains. Any future comparison should check
whether this threshold is compatible with $d_\mathrm{min}=0.02$ and the load
before searching for gains.

\paragraph{A blind spot of the specification.}
Controllers that are feasible with respect to the six constraints can
oscillate permanently. For an implementation-oriented journal, this is the
finding with the heaviest consequences: a mean-value constraint does not
guarantee the absence of a limit cycle. We do not modify the constraints after
observation; we recommend adding, in a future preregistered study, a
constraint on the oscillation amplitude over steady-state windows.

\paragraph{Fractional character.}
The best points found have $\lambda$ and $\mu$ equal or very close to~1. The
resulting controller is, in practice, of integer order. Nothing in our results
attributes a benefit to fractional order; the study lacks the direct-tuning,
equal-budget comparator that would allow such a judgement.

\section{Limitations and what has not been demonstrated}
\label{sec:limitations}

\begin{itemize}
\item \textbf{Model.} The campaign ran on a switched simulator written in
  Python. Cross-validation against Simulink covers only six candidates and
  the nominal plant; the snubber is assumed there ($R_s=10^5~\Omega$,
  $C_s=\infty$).
\item \textbf{Defect corrected and campaign rerun.} The stage-2
  feedforward first used the perturbed plant parameters. The campaign was
  rerun in full after correction (protocol v6.2); the conclusions are
  unchanged.
\item \textbf{Projection.} The FO-LQI $\to$ FO-PIDF identity is exact only
  in integer order, with the exact pole, on the unsaturated nominal plant; the
  pole projection alone introduces up to \ChiffreErreurProjPole{} error at the
  top of the band.
\item \textbf{Infeasibility not proven.} No cell is robustly feasible, but
  this does not prove that no robust controller exists: budgets are finite,
  and the uncertainty set is a declared choice.
\item \textbf{Uncertainty set.} The uncertain parameters and their ranges
  were chosen by us, as the brief does not specify them.
\item \textbf{Sealed test.} It was not played: no robustness claim is
  therefore made, not even as a proportion.
\item \textbf{Comparators.} The state-of-the-art references and the
  direct-tuning, equal-budget comparator are not implemented; no superiority
  can be attributed to fractional order.
\item \textbf{Post hoc diagnostics.} The voltage floor and the limit cycles
  were identified after observation; they motivate a future preregistered
  study but must not be read as confirmatory results.
\item \textbf{Metaheuristics.} Significant differences measure proximity to
  feasibility, not feasibility achieved, and their ranking depends on the arm.
\end{itemize}
